\documentclass[11pt]{article}

\usepackage[margin=1in]{geometry}
\usepackage{amsmath,amssymb}
\usepackage{booktabs}
\usepackage{enumitem}
\usepackage{graphicx}
\usepackage[hidelinks]{hyperref}
\usepackage{xcolor}
\usepackage{microtype}

\definecolor{caution}{RGB}{150,55,35}
\newcommand{\code}[1]{\texttt{#1}}
\newcommand{\caution}[1]{\textcolor{caution}{#1}}

\title{Clean Analysis of the June 2026 SAFOD\\
Accelerated Weight-Drop DAS Experiment}
\author{Analysis workflow documentation}
\date{1 August 2026}

\begin{document}
\maketitle

\begin{abstract}
This document describes a clean reanalysis of the June 2026 accelerated
weight-drop (AWD) experiment recorded by two distributed acoustic sensing
(DAS) installations at SAFOD.  The purpose of the workflow is to reconstruct
the experiment from authoritative timing information, preserve all valid AWD
bursts, identify coherent wave modes without assigning phase labels in
advance, and quantify the repeatability of each data-supported mode.  The
workflow is intentionally separate from previous exploratory notebooks.  In
particular, it does not use the locations of four conventional shots acquired
at other times, does not assume that fiber coordinate is calibrated vertical
depth, and does not call a coherent ridge a direct P wave or tube wave until
independent evidence supports that interpretation.
\end{abstract}

\section{Scientific motivation}

The experiment provides repeated, nominally identical controlled-source
observations on two different borehole DAS installations.  The AWD remained at
one surface location for the approximately 24-hour experiment, about 15~m
horizontally from the Nano wellhead.  Approximately 20 drops were performed in
each burst.  This design is well suited to three questions:

\begin{enumerate}[label=\arabic*.]
    \item Which coherent propagating modes are excited by the AWD source?
    \item How do those modes differ between the cemented Nano fiber and the
          wireline Deep fiber over a matched fiber-coordinate aperture?
    \item How stable are their moveout, waveform coherence, and amplitude from
          burst to burst?
\end{enumerate}

Velocity structure, attenuation, coupling, and time-lapse velocity change are
potential later products.  They are not treated as established results until
the observed modes and geometry have been validated.

\section{Data and working geometry}

\subsection{Nano fiber}

The Nano data are Sintela protobuf files in
\begin{quote}
\code{/oak/stanford/groups/ettore88/data/SAFOD/ActiveJune2026/Nano/}.
\end{quote}
The inventory contains 589 files.  Each file holds approximately 300~s of data
from 732 channels sampled at 1000~Hz.  Acquisition metadata and the existing
reader report a channel spacing of
\[
\Delta x_{\rm Nano}=1.26606202~\mathrm{m}
\]
and a gauge length of approximately 16.46~m.  The data returned by the current
reader are treated as strain-rate observations.

\subsection{Deep fiber}

The Deep data are HDF5 files from the simultaneously operating wireline DAS
array.  The existing paired product contains 3200 channels with a nominal
channel spacing of
\[
\Delta x_{\rm Deep}=2.0419~\mathrm{m}.
\]
The Deep fiber has a down-going and returning geometry.  Its fiber-coordinate
turnaround and true depth registration require separate validation.  The clean
mode comparison therefore uses only a matched 80--440~m interval in fiber
coordinate and does not claim that equal fiber coordinates correspond exactly
to equal physical depths.

\subsection{Source geometry}

The AWD source was fixed approximately 15~m horizontally from the Nano
wellhead.  If the fiber segment is vertical and its depth origin is known, the
source--receiver distance would be
\begin{equation}
r(z)=\sqrt{z^2+h^2}, \qquad h\simeq 15~\mathrm{m}.
\end{equation}
The correction relative to a zero-offset assumption is most important near
the surface.  The current workflow nevertheless plots distance along fiber,
not calibrated depth, because an absolute channel-to-depth registration has
not yet been established.

\subsection{Excluded conventional shots}

The file \code{SAFOD\_ActiveSrc\_ShotLocs\&Times\_20260617.csv} describes four
shots performed in the same general area at other times.  Those events are
not the repeating AWD experiment.  Their locations and times are excluded
from every step documented here.

\section{Why the analysis was restarted}

The recent exploratory notebook contains useful code and diagnostic products,
but it accumulated multiple evolving hypotheses and inconsistent epoch
conventions.  In particular, one branch inferred 48 epochs from the number of
five-minute files and truncated the GPS-defined burst list.  Other cells
handled a 49-row stack by discarding the first row.  Consequently, downstream
results could depend on which notebook cell had most recently defined the
working array.

The clean workflow adopts the following rules:

\begin{itemize}
    \item GPS times define drops and bursts; file count does not define epochs.
    \item Every excluded drop must have an explicit, machine-readable reason.
    \item Old arrays and figures are retained for reproducibility but are not
          silently reused as authoritative inputs.
    \item Moveout is detected before a geological or borehole-mode name is
          assigned.
    \item Fiber coordinate is not labeled as depth without calibration.
\end{itemize}

\section{Stage 1: authoritative AWD manifest}

The script \code{build\_manifest.py} reads the GPS timing catalog
\code{p26.cc9.txt}, indexes all Nano and Deep files by their acquisition start
times, and associates each drop with its containing file.  Adjacent GPS picks
are assigned to the same burst when their separation is no greater than 60~s.
Groups with fewer than five drops are treated as non-burst anomalies.

For every accepted drop, the manifest records
\begin{itemize}
    \item burst and within-burst identifiers;
    \item absolute UTC time and time since burst start;
    \item the timing catalog's maximum correlation and previous-pick interval;
    \item Nano and Deep file names and offsets within those files;
    \item nominal Nano, Deep, and paired availability.
\end{itemize}

This reconstruction finds
\begin{center}
\begin{tabular}{lr}
\toprule
Quantity & Count\\
\midrule
GPS entries within Nano acquisition coverage & 989\\
Qualifying AWD bursts & 49\\
Drops in qualifying bursts & 988\\
Drops with nominal Nano coverage & 988\\
Drops with nominal Deep coverage & 926\\
\bottomrule
\end{tabular}
\end{center}
The single GPS entry not included among the 988 drops belongs to a group that
does not meet the minimum burst-size definition.  Nominal Deep coverage is
absent for the final three bursts.  ``Nominal coverage'' only means that a file
spans the GPS time; the stack reconstruction additionally tests whether the
entire requested waveform window lies inside the decoded array.

The output is \code{awd\_manifest.csv}.  It is the authoritative bookkeeping
table for all later stages.

\section{Stage 2: canonical paired stack reconstruction}

The rebuild uses all Nano files rather than imposing the earlier artificial
24-hour filename boundary.  Drops are windowed using
\[
t\in[-0.5,3.0]~\mathrm{s}
\]
relative to their GPS times.  Both arrays are read at 1000~Hz and initially
filtered from 1 to 100~Hz.  A paired drop is retained for comparative analysis
only when a complete window can be extracted from both fibers.

For burst $b$ with $N_b$ retained drops, the stored epoch stack is
\begin{equation}
\bar{u}_b(x,t)=\frac{1}{N_b}\sum_{k=1}^{N_b}u_{bk}(x,t).
\end{equation}
The canonical product stores the Nano and full Deep stacks, $N_b$, burst start
times, channel spacings, sampling rate, and Deep channel bounds.  It is written
to
\begin{quote}
\code{canonical\_epoch\_stacks\_paired\_deep\_all.npz}.
\end{quote}
This new filename prevents accidental overwriting of the older stack product.

Stacking improves visibility of repeatable phases, but it suppresses
within-burst variability.  Therefore, stack-based mode identification is a
first robust visualization, not a substitute for later individual-drop QC.

\section{Stage 3: phase-neutral coherent-moveout scan}

\subsection{Matched apertures and frequency bands}

Nano and Deep are initially compared over the same 80--440~m fiber-coordinate
interval.  Three non-overlapping or minimally overlapping bands are analyzed:
\[
5\text{--}15~\mathrm{Hz},\qquad
15\text{--}30~\mathrm{Hz},\qquad
30\text{--}60~\mathrm{Hz}.
\]
The multiband design tests whether a ridge has stable kinematics or disperses
with frequency.  Such behavior is important for distinguishing body-wave and
guided-wave interpretations.

\subsection{Semblance definition}

For a trial unsigned apparent speed $v$ and intercept $t_0$, samples are
extracted along
\begin{equation}
t(x)=t_0+\frac{x-x_{\min}}{v}.
\end{equation}
The search covers 500--6000~m/s in 25~m/s increments and intercepts from
$-50$ to 300~ms in 2~ms increments.  A 40~ms waveform window is extracted at
each channel.  For channel samples $u_i(\tau)$, semblance is
\begin{equation}
S(v,t_0)=
\frac{\displaystyle\sum_{\tau}
      \left[\sum_{i=1}^{N}u_i(\tau)\right]^2}
     {\displaystyle N\sum_{\tau}\sum_{i=1}^{N}u_i^2(\tau)}.
\end{equation}
Perfectly aligned identical waveforms approach $S=1$, whereas incoherent
energy produces a much smaller value.  The analysis reports several distinct
speed peaks rather than only the global maximum, because multiple wave modes
may coexist.

The primary visualization, \code{mode\_scan.png}, contains six panels: three
bands for Nano and three bands for Deep.  Its numerical counterpart is
\code{mode\_scan.npz}.

\subsection{Interpretation discipline}

A semblance maximum is initially described only by its apparent speed,
intercept, bandwidth, and fiber.  A phase label requires additional evidence:
\begin{itemize}
    \item onset and polarity in an unnormalized record section;
    \item persistence across neighboring frequency bands;
    \item frequency-dependent speed or waveform dispersion;
    \item expected geometry after applying the 15~m source offset;
    \item consistency or contrast between the cemented and wireline fibers;
    \item stability across bursts and, ideally, individual drops;
    \item comparison with physically plausible formation and borehole-mode
          speeds.
\end{itemize}

\section{Stage 4: burst-to-burst repeatability}

The repeatability stage takes the strongest data-selected mode in each
fiber--band combination and performs a local search for every burst.  The local
speed search spans $\pm300$~m/s around the all-burst optimum, and the intercept
search spans $\pm20$~ms.  Each burst yields
\begin{itemize}
    \item best local apparent speed;
    \item best intercept;
    \item peak semblance;
    \item RMS amplitude along the aligned moveout;
    \item number of common drops contributing to that burst.
\end{itemize}

The output table \code{mode\_repeatability.csv} preserves every measurement.
The figure \code{mode\_repeatability.png} plots apparent speed by burst, colored
by semblance.  The median and median absolute deviation are reported because
they are less sensitive than a mean and standard deviation to anomalous bursts.

This test distinguishes a stable propagating ridge from a peak that appears
only after averaging the full day.  However, grid spacing and bandwidth limit
the precision.  Scatter smaller than one speed-grid increment must not be
interpreted as physical velocity precision.

\section{Products and visualization}

All clean products are located in
\begin{quote}
\code{/home/groups/ettore88/nberrios/safod\_das\_git/notebooks/awd\_clean/}.
\end{quote}

\begin{center}
\begin{tabular}{p{0.41\textwidth}p{0.48\textwidth}}
\toprule
File & Purpose\\
\midrule
\code{awd\_manifest.csv} & Authoritative drop, burst, timing, and file coverage table.\\
\code{canonical\_epoch\_stacks\_paired\_deep\_all.npz} & Corrected 49-burst paired stack product.\\
\code{mode\_scan.png} & Phase-neutral multiband Nano/Deep semblance visualization.\\
\code{mode\_scan.npz} & Numerical semblance surfaces and scan axes.\\
\code{mode\_repeatability.png} & Burst-to-burst apparent-speed and coherence stability.\\
\code{mode\_repeatability.csv} & Machine-readable per-burst mode measurements.\\
\bottomrule
\end{tabular}
\end{center}

The PNG files can be opened directly in JupyterLab.  The NPZ and CSV products
are intended for reproducible downstream analysis rather than visual inspection.

\section{Current limitations}

\begin{enumerate}
    \item \textbf{Fiber coordinate is not yet calibrated depth.} Geological
          depths must not be attached to anomalies until the channel origin,
          cable path, and Deep turnaround are verified.
    \item \textbf{The source offset is approximate.} The reported 15~m value is
          adequate for preliminary moveout work but should be replaced by a
          surveyed source--wellhead distance if available.
    \item \textbf{Cross-fiber amplitudes are not absolutely calibrated.}
          Interrogator response, gauge length, optical configuration, and
          strain versus strain-rate conventions must be reconciled before an
          absolute cemented/wireline amplitude ratio is interpreted.
    \item \textbf{Semblance can favor strong guided energy.} A high-coherence
          ridge is not automatically a direct body wave.
    \item \textbf{Burst stacks hide individual-source variability.} Final
          uncertainty analysis should return to single-drop windows.
    \item \textbf{The current speed scan is unsigned.} Reverse-propagating or
          returning-leg energy on the Deep fiber requires signed-slowness and
          hairpin-aware analysis.
    \item \textbf{A detection threshold has not yet been established by a
          null test.} Depth permutation, time scrambling, or pre-source scans
          should be used to determine which semblance peaks exceed chance
          coherence.
\end{enumerate}

\section{Criteria for the next scientific decision}

After the mode-scan and repeatability products are complete, the next analysis
will be chosen from the evidence:

\begin{itemize}
    \item A fast, weakly dispersive, early and repeatable ridge consistent on
          both arrays would motivate a direct-P travel-time and velocity study.
    \item A slower, dispersive, high-amplitude ridge that differs strongly
          between installations would motivate a borehole-guided-mode and
          coupling study.
    \item Stable kinematics but variable amplitude would motivate a source and
          coupling repeatability study.
    \item Poor per-burst stability despite an all-day semblance peak would
          indicate that the apparent mode is stack-dependent and unsuitable
          for time-lapse inference.
\end{itemize}

\section{What is deliberately not claimed}

At this stage, the analysis does not claim a calibrated interval-$V_P$ profile,
a detected Earth-tide velocity perturbation, a confirmed tube wave, a precise
maximum monitoring depth, or an absolute cemented-versus-wireline sensitivity
ratio.  Each could become a valid result, but only after phase, geometry,
calibration, and null-hypothesis tests support it.

\section{Reproducibility}

The clean workflow consists of the following scripts:
\begin{enumerate}
    \item \code{build\_manifest.py};
    \item \code{run\_canonical\_rebuild.py};
    \item \code{analyze\_modes.py};
    \item \code{measure\_repeatability.py}.
\end{enumerate}
The raw source directories are read only.  All new tables, arrays, logs, and
figures are written beneath \code{awd\_clean}.  Previous notebooks and figures
are preserved and are not overwritten.

\end{document}
